\chapter{局限、风险与实施计划}
\chapterlead{初赛提交在 10 月 18 日，决赛在 11 月底。剩下的时间，先补评测，再补功能。}

\section{当前的局限}

\paragraph{样本少。}真实数据只有五家公司，抽取评测只在三家年报上做过，研报质量只在两家公司上看到了版本间的显著提升。
三家年报是开发集，规则通道的版式处理就是在它们上面调出来的，结论外推到其他公司需要新的评测集来检验。

\paragraph{纠错核查缺少端到端评测。}注入实验跳过了断言抽取，只衡量核对器本身。真实研报中的指标别名、期间指代、作者的取整习惯，都还没有被系统地测过。

\paragraph{模型评审对结构性缺陷不敏感。}整节删除类缺陷几乎判不出；重复评审的总分极差 7～8 分，超过我们定的 5 分目标。

\paragraph{功能边界。}季报只走模型通道，规则通道不识别季报表头；扫描版 PDF 直接拒收；估值链不覆盖银行、保险等金融企业；
强周期公司缺少 P/B、重置成本等更合适的估值锚。

\paragraph{版本对比不是对照实验。}v0、v1、v2 用的是不同公司、不同稿件，流程改动和公司特性没有完全分开。

\section{风险与应对}

\begin{table}[H]
\centering
\caption{主要风险与应对}
\label{tab:pl-risk}
\small
\begin{tabularx}{\linewidth}{@{}>{\raggedright\arraybackslash}p{3.4cm} L L@{}}
\toprule
\hd{风险} & \hd{表现} & \hd{应对} \\
\midrule
新公司版式与开发集差异大 & 规则通道佐证率下降，冲突增多，估值被阻断 & 评测集覆盖 15 家左右、不同交易所与行业；冲突时退回分析师核对，不影响正确性，只增加耗时 \\
\rowband
会话异常中断 & 成本上升，耗时拉长（目前约两成会话退出码非 0） & 闸门按状态推进，断点续跑；现场演示准备好已跑通的工作区作为后备 \\
现场网络或模型服务不可用 & 决赛现场无法实时生成 & 计算引擎、纠错核查、评测全部可离线运行；研报流水线预先跑好，现场演示溯源与重算 \\
\rowband
评审偏好同族模型 & 分数系统性偏高 & 程序判的条目继续扩大；增加跨家族模型评审对照；人工抽查 \\
版权与合规 & 卖方研报、媒体报道的使用 & 只登记摘句与出处；评测用研报不入库；研报声明不构成投资建议 \\
\bottomrule
\end{tabularx}
\end{table}

\section{实施计划}

图~\ref{fig:pl-gantt} 是从现在到决赛的时间安排，表~\ref{tab:pl-tasks} 列出各项任务的交付物与验收标准。

\begin{figure}[H]
\centering
\resizebox{\linewidth}{!}{\begin{tikzpicture}[x=0.255cm,y=0.62cm]
  % 时间轴：第 0 天 = 10 月 7 日，至 11 月 30 日（第 54 天）
  \def\W{54}
  \fill[paper] (0,0.5) rectangle (\W,-9.6);
  \foreach \d/\l in {0/10.7,7/10.14,11/10.18,17/10.24,24/10.31,31/11.7,38/11.14,45/11.21,54/11.30}{
    \draw[rule] (\d,0.5) -- (\d,-9.6);
    \node[font=\tiny\sffamily,text=mute,above] at (\d,0.5) {\l};}
  \draw[seal,thick] (11,0.9) -- (11,-9.6); \node[font=\scriptsize\heiti\bfseries,text=seal,above] at (11,1.25) {初赛提交};
  \draw[seal,thick] (51,0.9) -- (51,-9.6); \node[font=\scriptsize\heiti\bfseries,text=seal,above] at (51,1.25) {决赛};
  \newcommand{\task}[5]{% 行号 起 止 颜色 名称
    \node[anchor=east,font=\scriptsize\heiti] at (-0.6,-#1) {#5};
    \fill[#4,rounded corners=1.5pt] (#2,-#1+0.28) rectangle (#3,-#1-0.28);}
  \task{0}{0}{11}{seal!80}{计划书与视频}
  \task{1}{0}{24}{ink!75}{评测集收集与标注}
  \task{2}{7}{31}{ink!55}{纠错端到端评测}
  \task{3}{11}{31}{inkl!70}{劣化算子与消融实验}
  \task{4}{17}{38}{gold!80}{季报表头识别}
  \task{5}{17}{38}{gold!60}{研报样例与写作改进}
  \task{6}{11}{45}{agentc!70}{网页前端与溯源演示}
  \task{7}{24}{45}{agentc!50}{终端界面实时刷新}
  \task{8}{38}{51}{seal!55}{决赛样例与答辩准备}
\end{tikzpicture}}
\caption{实施时间表（2026 年 10 月 7 日至 11 月 30 日；决赛具体日期以通知为准）}
\label{fig:pl-gantt}
\end{figure}

\begin{table}[H]
\centering
\caption{任务、交付物与验收标准}
\label{tab:pl-tasks}
\small
\begin{tabularx}{\linewidth}{@{}>{\raggedright\arraybackslash}p{2.9cm} L L@{}}
\toprule
\hd{任务} & \hd{交付物} & \hd{验收标准} \\
\midrule
计划书与视频 & 本计划书 PDF；不超过 5 分钟的演示视频 & 10 月 18 日前提交 \\
\rowband
评测集收集与标注 & 约 15 家公司的 2025 年报、2026 年一季报与半年报；8～10 家公司的卖方研报，人工埋错 150～200 处，另放约 20\% 容易误报的正确断言 & 每份文档登记来源与 SHA-256；埋错两人一组，一人埋、一人复核标签 \\
纠错端到端评测 & 真实研报上的召回率、误报率、类型准确率，按错误类型与漏判原因分解 & 结果可由一条命令复现；漏判逐条归因到抽取或核对 \\
\rowband
劣化算子与消融 & 跨全文删除并同步删除图表的劣化算子；资料搜集、复核会话、深度专题三项消融 & 每个配置不少于 3 家公司；报告与评审噪声（$\pm4$ 分）的比较 \\
季报表头识别 & 规则通道识别“本报告期”“年初至报告期末” & 季报纳入双通道交叉校验；现有测试与样例评测不回归 \\
\rowband
研报样例与写作 & 收集 5～10 篇优秀卖方深度报告，提炼分析逻辑写入技能说明（不做行业模板） & 新一批研报在同一评分标准下不低于 v2 \\
网页前端与溯源 & 网页端查看研报、真值表、估值工作簿；点击研报中的数字跳到年报对应页 & 现场可离线演示 \\
\rowband
终端界面 & 运行中实时刷新会话进度与闸门状态 & 能完整展示一次流水线运行 \\
决赛准备 & 现场样例数据、预跑好的工作区、答辩材料 & 断网条件下可演示全部离线功能 \\
\bottomrule
\end{tabularx}
\end{table}

\section{团队}

团队由四名同学组成，来自对外经济贸易大学和中国人民大学。表~\ref{tab:pl-team} 是成员与分工。

\begin{table}[H]
\centering
\caption{团队成员与分工}
\label{tab:pl-team}
\small
\begin{tabularx}{\linewidth}{@{}l l L@{}}
\toprule
\hd{成员} & \hd{学校} & \hd{主要负责} \\
\midrule
陈沫樵 & 对外经济贸易大学 & 项目负责人；系统架构、计算引擎与研报流水线 \\
\rowband
许冰洁 & 对外经济贸易大学 & \tbd{（待填）} \\
赵千淳 & 中国人民大学 & \tbd{（待填）} \\
\rowband
刘盈羽 & 对外经济贸易大学 & \tbd{（待填）} \\
\bottomrule
\end{tabularx}
\end{table}

评测集的收集与人工标注、人工质检复核、网页前端由团队成员分工承担；人工埋错按“一人埋、一人复核”两两配对，避免标签本身出错。
